\chapter*{Preface}
\addcontentsline{toc}{chapter}{Preface}

This book began as an essay. The essay was prompted by a long interview with the computer scientist Timnit Gebru, whose central objection to the prevailing public discussion of advanced artificial intelligence is that the opposition between salvation and extinction attributes agency to an artifact and thereby displaces accountability from the persons and institutions that build and deploy it. The objection is accepted here as a constraint on method. A proposal concerning powerful technology is incomplete unless it states who decides, who answers for the outcome, and how the outcome is checked, and the chapters that follow attempt to meet that standard for a conditional question: if systems of substantially greater scientific and engineering capability were to exist, what should humanity ask them to accomplish?

The essay stated principles. The book works out the programs to which those principles apply. It is organized in six parts. The first states the premises, sorts the problems that human communities already recognize by the kind of obstacle each presents, and fixes the terms in which accountability is assigned. The second concerns instruments, namely the use of advanced systems to generate proposals, to run experiments, to simulate and to plan, together with the discipline of verification without which generation is not knowledge. The third concerns institutions, and treats education, care, public services, and the architecture of authority within which any robotic embodiment would operate. The fourth and fifth parts treat terrestrial and orbital megascale programs, each under a common standard of specification, and the sixth treats extreme longevity, reliability over very long intervals, and the constraint that prevents speculative benefits to distant generations from overriding present obligations.

Each chapter closes with a section of mathematical formalization. The formalizations are of two kinds. Some state a structural claim precisely and prove it, so that the claim can be seen to depend on stated assumptions and no others. Others compute a quantity whose order of magnitude disciplines discussion, such as the energy required to lift water over a continental distance or the taper of a tether that must support its own weight. In neither case do the formalizations amount to engineering designs. A concept that survives the arithmetic of its own premises has passed a first and weak test, and the book says repeatedly that passing it licenses nothing beyond further and more demanding analysis.

Two cautions apply throughout. The programs described in the fourth and fifth parts vary widely in maturity, and the readiness classification introduced in Chapter 16 is used to keep them apart; a reader should not infer from their appearance in the same volume that they are equally plausible. Second, the book takes no position on the empirical forecasts, whether optimistic or catastrophic, on which the public debate turns. The conditional form of the central question is intended to make that agnosticism a feature of the argument and not an omission from it.
